\section{Retrieval Pipeline Results}

\subsection{Feature Extraction}
Features were extracted from 1,290 YOLO-segmented fungal colony images
across 31 strains (8 species) from the original prepared dataset. Each
segment was processed through 11 feature extractors:

\begin{itemize}[nosep]
  \item \textbf{Finetuned DL:} EfficientNetB1, ResNet50, MobileNetV2
  \item \textbf{Pretrained DL:} EfficientNetB1, ResNet50, MobileNetV2
  \item \textbf{Traditional:} HOG, Gabor, ColorHistogram, ColorHistogramHS
  \item \textbf{Hybrid:} ColorHistogramHS + ResNet50 (concatenated)
\end{itemize}

Features were indexed into Qdrant vector database with cosine distance.

\subsection{Retrieval Accuracy Sweep (K=7)}
A comprehensive sweep of 402 configurations evaluated all combinations of
extractors $\times$ media selection strategies $\times$ aggregation methods.

\textbf{Media Selection Strategies:}
\begin{itemize}[nosep]
  \item \textbf{E1 (Same):} Query and search within same environment
  \item \textbf{E2 (All):} Search across all available environments
  \item \textbf{E3-xxx:} Use only the specified medium
  \item \textbf{E4-xxx:} Exclude the specified medium
\end{itemize}

\textbf{Aggregation Strategies:} weighted, uni, freq\_strength, relative

The latest local rerun with updated finetuned weights and corrected feature-indexing pipeline achieved a best overall retrieval accuracy of \textbf{78.6\%} (33/42 predictions correct), with the top configuration coming from the pretrained EfficientNetB1 backbone under \texttt{freq\_strength}, E1, and $K=7$.

\begin{table}[H]
\centering
\caption{Top retrieval configurations from K=7 sweep}
\label{tab:retrieval_best}
\begin{tabular}{@{}lcccc@{}}
\toprule
Extractor & Media & Aggregation & Accuracy & Correct/Total \\
\midrule
EfficientNetB1 (pret.) & E1 & freq\_strength & 0.786 & 33/42 \\
MobileNetV2 (pret.) & E1 & freq\_strength & 0.738 & 31/42 \\
ResNet50 (pret.) & E1 & weighted & 0.714 & 30/42 \\
MobileNetV2 Finetuned & E1 & weighted & 0.714 & 30/42 \\
\bottomrule
\end{tabular}
\end{table}

\begin{figure}[H]
\centering
\includegraphics[width=\textwidth]{retrieval_top_configs.png}
\caption{Top 20 retrieval sweep configurations (K=7)}
\label{fig:retrieval_top}
\end{figure}

\begin{figure}[H]
\centering
\includegraphics[width=\textwidth]{retrieval_media_heatmap.png}
\caption{Best accuracy by extractor and media selection strategy}
\label{fig:retrieval_heatmap}
\end{figure}

\subsection{5-Fold Cross-Validation}
Strain-level 5-fold cross-validation was performed using fold-specific
EfficientNetB1 finetuned models trained on YOLO segments. Each fold
excludes one test strain per species, ensuring no strain-level data leakage.

200 configurations were evaluated (2 env $\times$ 4 aggregation $\times$ 5 K values 
$\times$ 5 folds). The best configuration achieved \textbf{82.9\% mean accuracy}:

\begin{table}[H]
\centering
\caption{Best cross-validation configurations}
\label{tab:cv_best}
\begin{tabular}{@{}lcccc@{}}
\toprule
Env & Aggregation & K & Mean Acc. & Std \\
\midrule
E2 & freq\_strength & 3  & 0.829 & 0.091 \\
E2 & relative      & 3  & 0.825 & 0.090 \\
E2 & weighted      & 3  & 0.825 & 0.090 \\
E2 & uni           & 3  & 0.825 & 0.090 \\
E1 & freq\_strength & 7  & 0.821 & 0.090 \\
\bottomrule
\end{tabular}
\end{table}

\begin{figure}[H]
\centering
\includegraphics[width=\textwidth]{cv_fold_configs.png}
\caption{Best CV configurations bar chart}
\label{fig:cv_configs}
\end{figure}

\begin{figure}[H]
\centering
\includegraphics[width=0.7\textwidth]{fold_variance_new.png}
\caption{Per-fold accuracy across all configurations}
\label{fig:fold_variance}
\end{figure}

\begin{figure}[H]
\centering
\includegraphics[width=0.9\textwidth]{training_curves_folds.png}
\caption{Fold-specific EfficientNetB1 training curves}
\label{fig:training_folds}
\end{figure}

\subsection{Segmentation Method Comparison}
YOLO and K-means segmentation methods were compared head-to-head using the
best retrieval configuration (ResNet50 Finetuned, freq\_strength).

\textbf{Result: Both methods achieve identical best accuracy of 78.6\%} (33/42) under the latest local rerun.
The segmentation choice still does not materially change the best species-level retrieval score on this dataset, although the winning backbone is now the pretrained EfficientNetB1 rather than the earlier fine-tuned ResNet50 claim.

\begin{figure}[H]
\centering
\includegraphics[width=\textwidth]{segmentation_comparison.png}
\caption{YOLO vs K-means segmentation: retrieval accuracy comparison}
\label{fig:seg_comp}
\end{figure}

\begin{table}[H]
\centering
\caption{Segmentation comparison by environment strategy}
\label{tab:seg_comp}
\begin{tabular}{@{}lccc@{}}
\toprule
Environment & YOLO & K-means & Difference \\
\midrule
E1 (same)        & 0.786 & 0.786 & 0.000 \\
E2 (all media)   & 0.619 & 0.619 & 0.000 \\
Weighted + E1 best & 0.714 & 0.714 & 0.000 \\
Weighted + E2 best & 0.667 & 0.667 & 0.000 \\
Best fine-tuned config & 0.714 & 0.595 & 0.119 \\
Best pretrained config & 0.786 & 0.786 & 0.000 \\
\bottomrule
\end{tabular}
\end{table}

\subsection{Unknown Species Threshold Detection}
Using the regenerated diverse retrieval results (342 grouped samples: 48 known + 294 unknown), 576 threshold formula-algorithm combinations were evaluated to detect unknown species via top-5 neighbor score aggregation.

The best regenerated formula achieved \textbf{F1 = 0.341} (Precision = 0.234,
Recall = 0.625) using the normalized gap between s0 and s2 with ROC-based
threshold optimization. The extreme class imbalance (48 known vs 294 unknown
samples) limits achievable F1.

\begin{table}[H]
\centering
\caption{Top threshold detection strategies}
\label{tab:threshold_best}
\begin{tabular}{@{}lccccc@{}}
\toprule
Formula & Algorithm & F1 & Precision & Recall & Threshold \\
\midrule
gap\_0\_2 & roc\_opt & 0.341 & 0.234 & 0.625 & 0.077 \\
gap0\_r2 & f1\_grid & 0.339 & 0.233 & 0.625 & 0.076 \\
exp\_decay\_20 & f1\_grid & 0.337 & 0.234 & 0.604 & 0.152 \\
\bottomrule
\end{tabular}
\end{table}

\begin{figure}[H]
\centering
\includegraphics[width=\textwidth]{threshold_top_strategies_new.png}
\caption{Top 10 threshold strategies for unknown species detection}
\label{fig:threshold_top}
\end{figure}

\begin{figure}[H]
\centering
\includegraphics[width=0.8\textwidth]{threshold_known_unknown_distribution.png}
\caption{Known vs unknown species distribution in full\_prepared dataset}
\label{fig:threshold_dist}
\end{figure}

\subsection{Summary of Key Results}
\begin{table}[H]
\centering
\caption{Overall experiment results summary}
\label{tab:summary}
\begin{tabular}{@{}lcc@{}}
\toprule
Experiment & Best Result & Note \\
\midrule
Retrieval K=7 sweep  & 0.786 acc & EfficientNetB1 pretrained, E1, freq\_strength \\
5-fold CV            & 0.829 mean & E2 freq\_strength k=3 \\
YOLO vs K-means      & 0.786 both & Identical best retrieval score \\
Unknown threshold    & 0.341 F1    & gap\_0\_2, roc\_opt \\
\bottomrule
\end{tabular}
\end{table}
